\begin{abstract}
\noindent Smartwatches record a single-lead electrocardiogram (ECG), but most diagnostic deep-learning models assume the 12-lead clinical ECG. We asked how much multi-label diagnostic ability survives when the input is reduced from 12 leads to wearable-like lead sets, and which diagnoses a reduced lead set could screen for. On PTB-XL (21,388 records, official patient-level folds) we trained baselines, three architectures with different inductive biases (ResNet1d, CNN-BiLSTM, ECG-ViT) and four training strategies (from scratch, supervised transfer, a lead-agnostic ResNet with lead masks, and self-supervised masked signal modelling) on six lead configurations, for five-class diagnosis and normal-versus-abnormal triage (102 training runs, three seeds). Reduced configurations were judged with a pre-specified criterion combining non-inferior AUROC (paired-bootstrap lower bound $\geq -0.02$) and sensitivity $\geq 0.80$ at 90\% specificity. All deep models reached 12-lead macro AUROC 0.915--0.925. Moving to lead I lowered macro AUROC by 0.090 (95\% CI 0.081--0.100), with myocardial infarction losing most ($\Delta$AUROC $-0.128$). No diagnosis met the criterion with lead I; ST/T change with the six chest leads was the only robust exception. Transfer learning and more labels did not close the gap, indicating an information bottleneck; one lead-agnostic model served all six configurations at a cost of 0.003--0.010 AUROC, and self-supervised pre-training improved a 10\%-label Transformer by 0.015 on lead I. Frozen triage models kept their discrimination on 34,815 external records from another hospital, but decision thresholds needed local calibration.
\end{abstract}
